\documentclass[12pt]{article}
\usepackage{lmodern}
\usepackage[T1]{fontenc}
\usepackage[dvips,letterpaper,margin=0.75in,bottom=0.75in]{geometry}
\usepackage{cancel}
\usepackage{graphicx}
\usepackage{braket}
\usepackage{latexsym,amssymb,amsmath}
\usepackage{pdfpages}
\usepackage{xcolor}
\usepackage{capt-of}
\usepackage{amsmath}
\usepackage{cite}


\begin{document}



%%%%%%%%%%%%%%%%%%%%%%%%%%%%%%%%%%%%%%%%%%%%%%%%%%%%%%%%%%%%%%%%
%%%%%%%%%%%%%%%%%%%%%%%%%%%%%%%%%%%%%%%%%%%%%%%%%%%%%%%%%%%%%%%%
\title{Proposal for Revising the Undergraduate Curriculum \\
Department of Physics and Astronomy}
%%%%%%%%%%%%%%%%%%%%%%%%%%%%%%%%%%%%%%%%%%%%%%%%%%%%%%%%%%%%%%%%
%%%%%%%%%%%%%%%%%%%%%%%%%%%%%%%%%%%%%%%%%%%%%%%%%%%%%%%%%%%%%%%%

\maketitle

\section*{Revision of Existing Majors and Minors}
{\it Please submit proposals for revision of an individual major or minor via email to Associate Director Mary Vasquez at \texttt{mpvasquez@ucdavis.edu}}.\\[4pt]

\noindent
Here we present a brief version of our proposal to revise the
undergraduate Physics BS major, including the Astrophysics
specialization.  We have also included a long write-up that describes
our complete plan for revising the curriculum, including Applied
Physics majors.  Much of that document is outside the scope of this
proposal, but it is included for reference of those interested in the
complete plan.  For convenience, the prompts from the guidance in
``Establishment, Revision, and Discontinuation of Undergraduate
Academic Degree Programs'' have been reproduced here, with matching
section numbers and prompt letters.

\section*{1. Approvals}

\noindent
{\bf Prompt a:} {\it Include all faculty votes and approval letters as required by PPM 200-25.}\\[4pt]

\noindent
This proposal has been developed over more than five years within the
Department of Physics and Astronomy, led by the undergraduate
curriculum committee.  The proposal was subjected to an extensive
vetting process that involved assigned department readers who were not
involved in formulating the initial plan.  {\bf The proposal was
unanimously supported by the department in a December 2021 vote.}

\section*{2. Background and Rationale}

\noindent
{\bf Prompt a:} {\it Explain all proposed revisions to the currently approved program.}\\[4pt]

\noindent
The updated BS physics major requirements include the following changes:
\begin{enumerate}
\item PHY 9HE is removed from the Honors Physics sequence, and the
  course will be discontinued.
\item PHY 110C is removed, and the course will be discontinued.
\item The math requirements have been updated to include additional
  new options available.
\item The computational physics course requirements have been updated.
  The previous requirement (PHY 102 or PHY 104B) has been replaced
  by the new required course PHY 45 (Computational Physics).  Students are also
  required to take three one-unit computational lab courses taught
  alongside core physics courses: PHY 110L, 112L, 115L.
\item The lab sequence option PHY 116ABC has been replaced by required courses PHY 80
  and new courses PHY 117 and 118.
\end{enumerate}
The updated Astrophysics emphasis requirements include the following changes:
\begin{enumerate}
\item PHY 9HE is removed from the Honors Physics sequence, and the
  course will be discontinued.
\item PHY 110C is no longer an elective, as the course will be discontinued.
\item The math requirements have been updated to include additional
  new options available.
\item The PHY 102 (Computational Laboratory in Physics) requirement is eliminated.  
\item The PHY 116A elective has been replaced with PHY 117 and 118 electives.
\item PHY 158 (Galaxy Formation) has been added as an option along
  with PHY 151,152,153, and 156.
\end{enumerate}

\noindent
{\bf Prompt b.} {\it Explain the rationale for why the proposed revisions are necessary.}\\[4pt]

\noindent
The changes proposed here address observed defects in our current program:
\begin{enumerate}
\item  Student workload over four years of study was imbalanced, with too few
  physics courses in the sophomore year, and too many in the junior
  year. 
\item Integration between transfer students and four-year students,
  who have somewhat different backgrounds, needed to be improved. Our
  recent departmental Climate Survey confirmed this by showing a huge
  satisfaction gap between these groups. 
\item The curriculum did not reflect the explosive growth in
  computational methods and their applications.
\item The highly successful new lab course PHY 80 introduced in a
  previous program revision overlapped significantly with lab courses
  PHY 116A and 116C.
\end{enumerate}

\noindent
Combined with already implemented schedule changes for other course
offerings, eliminating the PHY 9HE requirement allows students following
the honors physics sequence to start on other requirements starting in
the winter quarter of the their sophomore year.  This allows them to
balance their physics requirements more evenly throughout their four
years.  An additional benefit of the already implemented schedule changes,
including twice per year offerings of PHY 104A (Introduction to
Mathematical Methods in Physics), is that the course load for
transfer students is more balanced as well.\\

\noindent
Eliminating the PHY 110C requirement for BS majors allows us to add
additional required computational physics units while remaining under
the L\&S 110 unit cap on required coursework.  The content of PHY 110A
and 110B has already been modified.  PHY 104A now includes a review of
vector calculus that was formerly part of 110A.  PHY 110C is also
eliminated as an elective options for majors with an Astrophysics
emphasis.\\

\noindent
PHY 102 (Computational Lab in Physics) was technically a one-unit
course, but it was generally taught more like a 3-4 unit course.  Even
with the addition (in a previous program revision) of PHY 40, it is
not possible to cover computational physics in a meaningful way as a
single one unit course.  For BS physics majors, the proposal keeps the
PHY 40 requirement, eliminates PHY 102, and adds new required courses
(PHY 45, 110L, 112L, 115L).  The PHY 102 requirement is eliminated for
majors with an Astrophysics emphasis as well.  Astrophysics students
continue to take PHY 40 and encounter additional computational physics
in their upper division course work (e.g. PHY 151, 152, 153, 156, and
158).\\

\noindent
PHY 80 (Experimental Techniques) was added as required course for all
majors in a previous revision, primarily to ensure that students were
adequately prepared for the advanced lab courses PHY 122A and 122B.
The course has been quite successful, however, there is significant
overlap with 116ABC.  The courses have been reorganized so that the
content in 116ABC is now contained in PHY 80, 117, and 118.\\

\noindent
We have been approving students requests to substitute newly available
math courses (MAT 27A, 67, 27B) where appropriate for our current
requirements (MAT 22A and 22B).  This proposal makes these
equivalencies explicit.\\

\noindent
Previously, majors with an Astrophysics emphasis were required to take
PHY 151, 152, 153, and 156.  In this proposal, students are given an
additional option, as they choose four courses from that list of four
plus the new course PHY 158 (Galaxy Formation).\\

\noindent
{\bf Prompt c.} {\it Describe the timing of and implementation plan for the revisions. Please be sure to address all impacts to students and catalog rights, including:}\\
{\bf i.} {\it How many existing students in the major will be affected by the proposed change?}\\
{\bf ii.} {\it Are any mitigation measures needed to prevent delays in time to degree for students already enrolled in the major or minor?}\\

\noindent
All of our approximately 150 physics BS majors will be affected by the
proposed changes.  However, all of the new courses have been taught
for several years now, and students have been advised to take them.
The necessary changes to course content and prerequisites in existing
courses have already been completed and approved in IPA.  We have also
already made the required adjustments to course scheduling.  Our
undergraduate advisor confirms that all current students will be able
to graduate under the new requirements without delaying their
graduation date.  Approving the new major requirements will clarify
the situation for our students, who have been hearing about these
changes for years, and so making the change as soon as possible would
be best for everyone.

There is a clear equivalency of requirements that will allow us to
handle special cases as needed:
\begin{center}
\begin{tabular}{ll}
Old Requirement & New Requirement \\
\hline
PHY 9HE & Dropped \\
PHY 110C & Dropped \\    
PHY 102 or PHY 104B & PHY 45, 110L, 112L, 115L \\  
PHY 116A, 116B, 116C & PHY 80, 117, 118 \\  
\end{tabular}
\end{center}
We are also willing to waive the one-unit lab course requirements
during the transition on a case-by-case basis, in case we find that it
is necessary.  However, we suspect that it will not be necessary.\\

\noindent
{\bf Prompt d.} {\it Each proposal should contain revisions for only
  one major or minor program. Occasionally, multiple related proposals
  need to be approved in sequence to avoid disruption to
  enrollments. Please indicate this in cover letters, but submit
  separate proposals.}\\

\noindent
This proposal is focused exclusively on BS Physics majors.  We expect
to submit updates to the AB Physics and Applied Physics majors
separately, but these can be approved independently.

\section*{3. Consultation with Affected Units (e.g., Departments, Programs, Centers)}

{\bf Prompt a.} {\it Describe overlap or conflict with existing UC Davis majors or minors.}\\

\noindent
The closest major to a BS Physics offered by another deparment is the
BS in Applied Mathematics.  Students in that major might select PHY 9A
and 9B plus one enrichment course that could included an upper-level
physics requirement such as 105A or 104A.  The focus is much more
heavily on mathematics, and overall this is not a significant overlap.
In terms of upper division course work, this is four units in physics
compared to 59 units (minimum) for our majors.\\

\noindent
{\bf Prompt b.} {\it You must consult with affected units and include letters from such units stating their feedback. Proposals may still be submitted even if affected units are not supportive. The Davis Division will only review proposals that include:}\\
{\bf i.} {\it Feedback from affected units; or}\\
{\bf ii.} {\it If affected units choose not to provide feedback, documentation of reasonable attempts to receive this feedback.}\\

\noindent
We have included letters of support from both Department of
Mathematics and the Department of Earth and Planetary Sciences.  In
the case of Mathematics, we are simply implementing changes to include
additional new variants, at their request.  In the case of Earth and
Planetary Sciences, we are keeping the same course as an elective, but
updating the name , also at their request.

\section*{4. Curricular Structure}

\noindent
{\bf Prompt a.} {\it Revisions to the curricular structure should be submitted in track changes (deletion and addition of courses and changes to total required units should be clear).}\\

\noindent
We have submitted the updated curricular stucture in track changes.\\

\noindent
{\bf Prompt b.}{\it Provide letters of support from units whose courses are added to the program in this revision.}\\
{\it i. From Chairs/Directors: Confirm courses have capacity to accommodate frequency of offerings and number of projected students; comment on any other levels of unit/program support required.}\\

\noindent
We have included letters of support from both Department of Mathematics and the Department of Earth and Planetary Sciences.\\

\noindent
{\bf Prompt c.}
{\it Note that as per COCI policy, UGC acknowledges all in-person, virtual (V), and hybrid (Y) versions of a course number to be interchangeable with respect to major and prerequisite requirements. In such cases where differentiation is desired (one version accepted toward the degree and another not), the program should contact COCI about required adjustments to course numbering.}\\

\noindent
We have no need to make such a request at this time.\\

\noindent
{\bf Prompt d.}
{\it Consider the impact of prerequisites to required courses that are not explicitly listed among major requirements (“hidden prerequisites”) on time to degree.}\\

\noindent
We have been careful to remove any hidden prequisites.  We have added explicit requirements
for prerequisites as needed for all required courses.\\

\noindent
MAT 22A has a concurrent prerequisite for MAT 22AL, but this has not
been enforced, and from our recent discussions with the Department of
Mathematics we understand that this prequisite is likely to be
eliminated soon.  We therefore prefer to leave this out of our
requirements (as in the currently approved requirements).\\

\end{document}

